\documentclass[11pt,a4paper]{article}
\usepackage[margin=2cm]{geometry}
\usepackage{titlesec}
% Preámbulo común de los documentos (article). Compatible con Overleaf (pdfLaTeX, sin -shell-escape).
\usepackage[utf8]{inputenc}
\usepackage[T1]{fontenc}
\usepackage[spanish,es-noshorthands]{babel}
\usepackage{lmodern}
\usepackage{microtype}
\usepackage{amsmath,amssymb}
\usepackage{booktabs}
\usepackage{array}
\usepackage{tabularx}
\usepackage{xcolor}
\usepackage{listings}
\usepackage{enumitem}
\usepackage{caption}
\usepackage[hidelinks]{hyperref}

\captionsetup{font=small,labelfont=bf}
\setlist{itemsep=1pt,topsep=3pt}

\lstset{
  basicstyle=\ttfamily\footnotesize,
  columns=fullflexible,
  keepspaces=true,
  frame=single,
  framerule=0.3pt,
  rulecolor=\color{gray!60},
  backgroundcolor=\color{gray!6},
  breaklines=true,
  inputencoding=utf8,
}

\newcommand{\rel}[1]{\textbf{#1}}
\newcommand{\code}[1]{\texttt{#1}}

\titlespacing*{\section}{0pt}{8pt}{3pt}
\titlespacing*{\subsection}{0pt}{5pt}{2pt}
\setlength{\parskip}{2pt}

\title{\vspace{-1.2cm}Listwise Learning to Rank con AdaRank\\[1pt]
\large Construcción de un training set de relevancia para consultas LOINC}
\author{Grupo: [Integrante 1], [Integrante 2], [Integrante 3]\\
\small Machine Learned Ranking --- MUIA, Universidad Politécnica de Madrid}
\date{\vspace{-0.4cm}25 de noviembre de 2026}

\begin{document}
\maketitle

\section{Introducción y enfoque}

El objetivo es ordenar, para cada consulta clínica, una lista de pruebas de laboratorio codificadas en LOINC
según su relevancia. En \emph{Learning to Rank} hay tres familias. Los métodos \emph{pointwise} predicen la relevancia de cada documento por separado.
Los \emph{pairwise} aprenden a ordenar pares de documentos. Los \emph{listwise} optimizan directamente una medida de calidad de la lista
completa de cada consulta. Hemos elegido el enfoque \textbf{listwise}, porque lo que importa al usuario es el orden
de la lista y no la puntuación aislada de cada documento. Como algoritmo usamos \textbf{AdaRank} \cite{xu2007adarank}.

\paragraph{AdaRank.} Es un algoritmo de \emph{boosting} que construye un ranker lineal
$f_T(\mathbf{x})=\sum_{t=1}^{T}\alpha_t h_t(\mathbf{x})$, donde cada ranker débil $h_t$ es una feature.
Mantiene una distribución de pesos $P_t$ sobre las \emph{consultas}, que empieza uniforme. En cada ronda:
\begin{enumerate}[nosep]
\item Elige la feature con mayor $\sum_i P_t(i)\,E(\pi(q_i,h),\mathbf{y}_i)$, donde $E$ es la medida de ranking (aquí NDCG@10).
\item Le da el peso $\alpha_t=\tfrac12\ln\frac{\sum_i P_t(i)\{1+E(\pi(q_i,h_t),\mathbf{y}_i)\}}{\sum_i P_t(i)\{1-E(\pi(q_i,h_t),\mathbf{y}_i)\}}$.
\item Actualiza $P_{t+1}(i)\propto\exp(-E(\pi(q_i,f_t),\mathbf{y}_i))$, de modo que las consultas peor ordenadas ganan peso.
\end{enumerate}
La unidad de optimización es, por tanto, la lista de documentos de cada consulta. Usamos como métrica
$\mathrm{NDCG@}k=\mathrm{DCG@}k/\mathrm{IDCG@}k$, con $\mathrm{DCG@}k=\sum_{r=1}^{k}(2^{y_r}-1)/\log_2(r+1)$ \cite{jarvelin2002}.

\section{Dataset y etiquetas de relevancia}

El libro \code{loinc\_dataset-v2.xlsx} tiene tres hojas, una por consulta: \emph{glucose in blood} (qid 1),
\emph{bilirubin in plasma} (qid 2) y \emph{white blood cells count} (qid 3). Cada hoja contiene los mismos 67 códigos
LOINC, en distinto orden, con las columnas \code{loinc\_num}, \code{long\_common\_name}, \code{component},
\code{system} y \code{property}. Esto da 201 pares consulta--documento. El Excel no incluye etiquetas, y el orden de las filas
no se usa como tal.

\paragraph{Escala 0--3.} Cada par se juzga en tres dimensiones:
\begin{itemize}[nosep]
\item \textbf{analito} (\code{component}): exacto, relacionado (fracción o subpoblación) o ninguno;
\item \textbf{muestra} (\code{system}): exacta, compatible pero distinta, o incompatible;
\item \textbf{medida/condición} (\code{property} y sufijo \code{\^{}}): adecuada o desviada.
\end{itemize}
La regla es determinista:
\begin{itemize}[nosep]
\item un analito distinto da \rel{0};
\item si no, se parte de \rel{3} (analito exacto) o \rel{2} (relacionado);
\item se resta 1 si la muestra es compatible pero distinta y 2 si es incompatible;
\item se resta 1 si la medida está desviada;
\item el mínimo es \rel{1} si se comparte analito.
\end{itemize}
Las decisiones clave fueron:
\begin{itemize}[nosep]
\item \code{Ser/Plas} es compatible pero distinta de \code{Bld}, así que la glucosa en suero/plasma recibe \rel{2};
\item WBC = leucocitos, y un \emph{porcentaje} de linfocitos no es un recuento (\rel{1});
\item la glucosa en ayunas o post-carga es una prueba funcional (medida desviada);
\item coincidencias léxicas como \emph{Blood product type} reciben \rel{0}.
\end{itemize}
La regla está implementada en \code{src/queries.py} y cada par lleva su justificación textual en
\code{data/relevance\_labels.csv}. La guía completa está en el anexo \code{labeling\_guidelines.pdf}. Respecto a la
versión inicial (0--2), se corrigió la errata \emph{bilirrubin} de la consulta y cinco juicios incoherentes.

\begin{table}[h]
\centering\footnotesize
% Generado por src/export_tables.py — no editar a mano
\begin{tabular}{llp{6.2cm}llc}
\toprule
Consulta & LOINC & long\_common\_name & system & property & Rel. \\
\midrule
glucose in blood & 74774-1 & Glucose [Mass/volume] in Serum, Plasma or Blood & Ser/Plas/Bld & MCnc & 3 \\
 & 14749-6 & Glucose [Moles/volume] in Serum or Plasma & Ser/Plas & SCnc & 2 \\
 & 15076-3 & Glucose [Moles/volume] in Urine & Urine & SCnc & 1 \\
 & 14747-0 & Glucose [Moles/volume] in Pleural fluid & Plr fld & SCnc & 1 \\
 & 14764-5 & Glucose [Moles/volume] in Serum or Plasma --3 hours post 100 g glucose PO & Ser/Plas & SCnc & 1 \\
 & 35184-1 & Fasting glucose [Mass or Moles/volume] in Serum or Plasma & Ser/Plas & MSCnc & 1 \\
\midrule
bilirubin in plasma & 1975-2 & Bilirubin.total [Mass/volume] in Serum or Plasma & Ser/Plas & MCnc & 3 \\
 & 1968-7 & Bilirubin.direct [Mass/volume] in Serum or Plasma & Ser/Plas & MCnc & 2 \\
 & 1971-1 & Bilirubin.indirect [Mass/volume] in Serum or Plasma & Ser/Plas & MCnc & 2 \\
 & 35192-4 & Bilirubin.indirect [Mass or Moles/volume] in Serum or Plasma & Ser/Plas & MSCnc & 2 \\
 & 14423-8 & Bilirubin.total [Mass/volume] in Synovial fluid & Synv fld & MCnc & 1 \\
 & 33870-7 & Bilirubin.total [Presence] in Unspecified specimen & XXX & PrThr & 1 \\
\midrule
white blood cells count & 26464-8 & Leukocytes [\#/volume] in Blood & Bld & NCnc & 3 \\
 & 26474-7 & Lymphocytes [\#/volume] in Blood & Bld & NCnc & 2 \\
 & 26484-6 & Monocytes [\#/volume] in Blood & Bld & NCnc & 2 \\
 & 26478-8 & Lymphocytes/100 leukocytes in Blood & Bld & NFr & 1 \\
\bottomrule
\end{tabular}

\caption{Documentos con relevancia $>0$. Los otros 61--63 documentos de cada consulta son \rel{0}.}
\end{table}

\section{Features y formato del training set}

Cada par se describe con cinco features que dependen de la consulta y se calculan \emph{sin} mirar la etiqueta
(detalle en \code{feature\_dictionary.pdf}):
\begin{itemize}[nosep]
\item $f_1$: coseno TF-IDF entre la consulta y \code{long\_common\_name};
\item $f_2$: BM25 normalizado por consulta \cite{robertson2009};
\item $f_3$: solapamiento de términos con sinónimos (\emph{white cells}$\to$\emph{leukocytes}, \emph{blood}$\to$\emph{bld});
\item $f_4$: el analito de la consulta aparece en \code{component} (0/1);
\item $f_5$: el término de muestra aparece en \code{system} (0/1).
\end{itemize}
Las codificaciones \code{LabelEncoder} del trabajo inicial se eliminaron. Eran índices alfabéticos, iguales para cualquier consulta,
y no tienen sentido como rankers débiles. El training set se exporta en formato LETOR (\code{data/train.txt}): 201 líneas
agrupadas en 67 por \code{qid}, con el código LOINC en comentario. Un script comprueba los recuentos, que el conjunto de códigos es el mismo en cada consulta,
que las etiquetas están en $\{0,\dots,3\}$, que las features son finitas y 16 casos de referencia.

% Generado por src/export_tables.py — no editar a mano
\begin{lstlisting}
3 qid:1 1:0.750959 2:1 3:1.01765 4:1 5:1 6:1 # docid = 74774-1
2 qid:1 1:0.499538 2:0.694848 3:0.465047 4:1 5:0 6:1 # docid = 14749-6
1 qid:1 1:0.428328 2:0.75537 3:0.5 4:1 5:0 6:1 # docid = 15076-3
3 qid:2 1:0.653337 2:1 3:1.03668 4:1 5:1 6:1 # docid = 1975-2
3 qid:3 1:0.467694 2:0.965348 3:1.01449 4:0.333333 5:1 6:1 # docid = 26464-8
1 qid:3 1:0.308181 2:0.881265 3:0.761091 4:0.333333 5:1 6:0 # docid = 26478-8
0 qid:1 1:0.237446 2:0.458793 3:0.491787 4:0 5:0 6:0 # docid = 933-2
\end{lstlisting}


\section{Verificación, resultados y limitaciones}

Con solo tres consultas no hay partición train/test fiable. Hacemos \textbf{leave-one-query-out}: cada fold entrena con dos consultas y evalúa en la tercera, sin
repartir documentos de una misma consulta entre train y test. Comparamos con cada feature usada sola como ranker.
Como ampliación opcional, añadimos tres consultas (\emph{calcium in serum}, \emph{cholesterol in plasma},
\emph{chloride in blood}) etiquetadas con la misma regla sobre los mismos 67 documentos.

\begin{table}[h]
\centering\small
% Generado por src/export_tables.py — no editar a mano
\begin{tabular}{lrrrr}
\toprule
Sistema & NDCG@1 & NDCG@3 & NDCG@5 & NDCG@10 \\
\midrule
\multicolumn{5}{l}{\emph{3 consultas oficiales (3 folds)}} \\
\textbf{AdaRank} & 0.667 & 0.772 & 0.796 & 0.797 \\
solo $f_1$ (coseno TF-IDF) & 0.667 & 0.856 & 0.889 & 0.889 \\
solo $f_2$ (BM25) & 0.667 & 0.764 & 0.829 & 0.863 \\
solo $f_3$ (solap. sinónimos) & 1.000 & 0.875 & 0.885 & 0.877 \\
solo $f_4$ (analito en component) & 0.714 & 0.824 & 0.821 & 0.824 \\
solo $f_5$ (muestra en system) & 0.333 & 0.366 & 0.475 & 0.467 \\
\midrule
\multicolumn{5}{l}{\emph{6 consultas, ampliado$^\dagger$ (6 folds)}} \\
\textbf{AdaRank} & 0.833 & 0.928 & 0.944 & 0.944 \\
solo $f_1$ (coseno TF-IDF) & 0.833 & 0.928 & 0.944 & 0.944 \\
solo $f_2$ (BM25) & 0.833 & 0.882 & 0.914 & 0.931 \\
solo $f_3$ (solap. sinónimos) & 0.905 & 0.878 & 0.878 & 0.893 \\
solo $f_4$ (analito en component) & 0.619 & 0.820 & 0.832 & 0.834 \\
solo $f_5$ (muestra en system) & 0.167 & 0.183 & 0.237 & 0.234 \\
\bottomrule
\end{tabular}

\caption{NDCG medio en leave-one-query-out. $^\dagger$Incluye las consultas adicionales.}
\end{table}

\paragraph{Resultados.} El modelo entrenado con las tres consultas ($\alpha_{f_1}{=}2.82$, $\alpha_{f_3}{=}1.41$) ordena
perfectamente su propio conjunto: todos los documentos relevantes quedan arriba, en orden de etiqueta. También ordena bien una consulta no vista
(\emph{cholesterol in plasma}, NDCG@10 = 1.0). En leave-one-query-out, en cambio, AdaRank \emph{no} supera a la mejor feature
individual. Con dos consultas de entrenamiento elige la feature que mejor ordena esas dos, y esa feature no tiene por qué servir para la tercera.
Por ejemplo, $f_1$ falla en \emph{white blood cells count} porque la consulta no comparte vocabulario con \emph{Leukocytes}, mientras que $f_3$
(con sinónimos) acierta. Con seis consultas, AdaRank iguala al mejor baseline (NDCG@10 = 0.944). Esto apunta a que el
cuello de botella es el número de consultas, no el algoritmo.

\paragraph{Limitaciones.}
\begin{itemize}[nosep]
\item Tres consultas son muy pocas para estimar la generalización; los resultados son demostrativos.
\item Las etiquetas son muy dispersas: 4--6 documentos relevantes de 67.
\item $f_4$ y $f_5$ se inspiran en las mismas columnas que la regla de etiquetado, aunque son más gruesas.
\item Las features binarias producen empates, y NDCG resuelve los empates según el orden de entrada.
\item En la implementación de AdaRank \cite{adarankgit} corregimos una división por cero: si un ranker débil es perfecto en todas las
  consultas, $\alpha=\infty$ y las predicciones salen NaN.
\end{itemize}

\section{Contribuciones}

\begin{itemize}[nosep]
\item \textbf{[Integrante 1]}: implementación de AdaRank y métricas, primera versión del dataset (0--2) y notebook inicial.
\item \textbf{[Integrante 2]}: escala de relevancia 0--3 y guía de etiquetado, features, exportación LETOR,
  evaluación leave-one-query-out y documentación.
\item \textbf{[Integrante 3]}: [completar].
\end{itemize}

\begin{thebibliography}{9}\small\setlength{\itemsep}{0pt}
\bibitem{xu2007adarank} J.~Xu y H.~Li. AdaRank: a boosting algorithm for information retrieval. \emph{SIGIR}, 2007.
\bibitem{jarvelin2002} K.~Järvelin y J.~Kekäläinen. Cumulated gain-based evaluation of IR techniques. \emph{ACM TOIS}, 20(4), 2002.
\bibitem{robertson2009} S.~Robertson y H.~Zaragoza. The probabilistic relevance framework: BM25 and beyond. \emph{FnTIR}, 3(4), 2009.
\bibitem{adarankgit} R.-C.~Chen. AdaRank (implementación en Python). \url{https://github.com/rueycheng/AdaRank}.
\end{thebibliography}

\end{document}
